\scanpage{3}

\uline{Polyèdres combinatoires}

\diagram{$\left.\begin{array}{ccccc} S&&A&&F\\ &R&&R'&\end{array}\right|\ \uline{\text{finis}}$}

\diagram{\begin{tikzcd}[ampersand replacement=\&,row sep=small,column sep=small]
R\subset S\times A\arrow[dr,dash]\&{}\&R'\subset A\times F\arrow[dl,dash]\\
\&\text{\textquotedblleft relations d'incidences\textquotedblright}\&
\end{tikzcd}}

\[
\Rightarrow\ \Phi\ \deleted{=}\ S\amalg A\amalg F
\quad\text{ens. ordonné}
\]

\[
 R''\subset S\times F,\quad
 \underset{\substack{\Vert\\R\circ R'}}{R''}=R\times_A R'
\]

\margin{\deleted{\unsure{poly}} ``géométrie d'incidences'' de dim 2}
% margin note written sideways

1) $\forall\,a\in A,\ \operatorname{card}R(a)=\operatorname{card}F(a)=2$

2) $\forall\,s\in S,\ f\in F$ avec $s<f$,

\[
 \operatorname{card}(\{a\in A\mid s<a<f)=2
\]

3) $\forall\, f\in F$, l'ens ordonné $\Phi_{<f}=\{x\in\Phi\mid x<f\}$ (qui est un \unsure{ordon|ordun|contour}) est \uline{connexe} \inserted{$\ne\varnothing$} i.e. est un polygone combinatoire (à $q(f)$ $\geq 2$ côtés)

3$'$) $\forall s\in S$, l'ens ordonné $\Phi_{>s}=\{x\in\Phi\mid x>s\}$ (``\quad '') est connexe $\ne\varnothing$ i.e., est un polygone combinatoire (à pfs $\geq 2$ côtés)

\margin{\deleted{\unsure{on exige}} \unsure{sup$^{t}$} ``polyèdre combinatoire''}
% slanted

4) Dans $\Phi$ les sup majorés existent ($\Longleftrightarrow$ les inf minorés existent)

\inserted{i.e. $x,y\in\Phi,\ \substack{x\nleq y\\ y\nleq x}\mid x,y<f\in F\Rightarrow\operatorname{Sup}(x,y)$ existe}

i.e. $\left\{\begin{aligned}
 &\text{a) Toute arête $a$ est déterminée par l'ens $\{s,s'\}$ de ses extrémités, et tte}\\
 &\text{face qui majore celles-ci majore l'arête.}\\
 &\text{b) Toute face $f$ est connue quand on connait une arête $a$ de $f$ et}\\
 &\text{un sommet de $f$ non incident, et si on connait deux}\\
 &\text{sommets de $f$ non incidents à une même arête.}
\end{aligned}\right.$

\uline{NB} Indépendamment de \uline{4)}, si $s,s'\in S,\ s\ne s',\ s,s'<a\in A$, \deleted{\unsure{alors|alors si}}

1°) $\operatorname{Sup}(s,s')$ existe sss $a=\operatorname{Sup}(s,s')$ -- \deleted{il \inserted{suff.} \unsure{faut} \inserted{sss} que $\operatorname{Sup}(s,s')$ existe, il f. et s. que} a) \deleted{On} soit seule arête qui soit incidente à $s,s'$

b) Toute face incidente à la fois à $s,s'$ l'est à $a$).

\deleted{2°) $\operatorname{Sup}(s,s')$ existe \unsure{sss}}

2°) Si $s,s'\in S,\ s,s'<f\in F$, \deleted{$s\ne s'$}, non incidents à une \deleted{\unsure{seule|même}} arête, alors $\operatorname{Sup}(s,s')$ existe \deleted{$\Longleftrightarrow$} $f=\operatorname{Sup}(s,s')$ \deleted{$\Longleftrightarrow$} $f$ l'unique face qui contient \inserted{$s,s'$}

3°) Si $s\in S,\ a\in A$ non incidents, et $s,a<f\in F$, alors $\operatorname{Sup}(s,a)$ existe sss $f=\operatorname{Sup}(s,a)$, \deleted{et alors \unsure{on}} i.e. sss $f$ est l'unique face majorant $s$ et $a$

3°) Si $a\in A,\ a'\in A,\ a\ne a',\ a,a'<f\in F$, alors $\operatorname{Sup}(a,a')$ existe sss $f=\operatorname{Sup}(a,a')$ i.e. sss $f$ est l'unique face majorant $a,a'$.

(et relations duales)

\diagram{\begin{tikzcd}[ampersand replacement=\&,column sep=small]
{}\arrow[d,dash,bend right=36,start anchor=west,end anchor=west]\&{}\\
{}\arrow[r,Leftrightarrow]\&{}\arrow[u,dash,bend right=24]
\end{tikzcd}}

$a_2)$ $a\in A$ connue quand on connait l'ens de ses sommets, i.e. $A\hookrightarrow\mathfrak P_2(S)$

$b_2)$ Si $f,f'\in F,\ f\ne f'$, alors l'ens des facettes inc. à $f,f'$ est soit vide, soit réduit à un sommet, soit formé d'une arête et de ses 2 sommets

on aura

\margin{\uline{On aura} $a'_1)$ $A\hookrightarrow\mathfrak P_2(F)$ $b'_1)$ Si $s,s'\in S,\ s\ne s'$, l'ens des facettes inc. à $s,s'$ est soit vide, soit réduit à une seule face, soit formé d'une arête et des deux faces incidentes à celles-ci.}
% slanted

4$'$) (\unsure{équiv.} à 4) a$'$) Toute arête $a$ est connue quand on connait l'ens $\{f,f'\}$ des deux faces incidentes, et tt sommet qui est \uline{incident} \deleted{\unsure{à}} à celles-ci est incident à l'arête.

b$'$) Un sommet $s$ est connu, quand on connait une arête $a$ et une face $f$ \deleted{incident\unsure{e|es} de $s$}, non incidentes

Pour un \uline{polyèdre combinatoire}, on a

\uline{Corollaire} 1°) Si à toute facette on associe l'ens des sommets qu'elle majore, on trouve une application

\[
 \Phi\longrightarrow\mathfrak P(S)
\]

qui est injective, et induit un iso. \inserted{d'ens. ordonnés} de $\Phi$ avec

% archive sheet number: 2

